% Job posting: https://www.linkedin.com/jobs/view/4450566062
\documentclass[11pt]{article}
\usepackage[left=1in,top=1in,right=1in,bottom=1in]{geometry}
\usepackage[hidelinks]{hyperref}
\usepackage{setspace}
\usepackage[T1]{fontenc}
\usepackage{newtxtext,newtxmath}
\ifdefined\pdfgentounicode
  \input{glyphtounicode}
  \pdfgentounicode=1
\fi
\setstretch{1.1}
\pagenumbering{gobble}
\begin{document}
\pagestyle{empty}

\noindent Nishal Stanislaus Silva, PhD \\
Toronto, ON, Canada \\
nishal.silva@hotmail.com \,|\, +1 613 200 2030 \\
\href{https://nishal.xyz}{nishal.xyz} \,|\, \href{https://www.linkedin.com/in/nishal-silva}{linkedin.com/in/nishal-silva} \,|\, \href{https://github.com/ns2max}{github.com/ns2max}

\vspace{1em}

\noindent Dear Hiring Manager,

\vspace{0.5em}

\noindent I am writing to apply for the Staff Engineer, Computer Vision role at Stryker in Burnaby, BC. Early in my career at MAS Holdings, South Asia's largest apparel manufacturer, I built and deployed a COVID-19 remote-vitals computer-vision device across hospital wards -- digitizing patient vitals under real clinical adoption stakes, a deployment recognized by the Government Medical Officers' Association (GMOA) of Sri Lanka. That project remains my most direct experience of what it means to ship a CV system into a clinical setting, where correctness and reliability carry real consequences. It sat alongside a broader body of production CV work at MAS: OpenCV-based multilingual care-label QC with an explainable defect overlay integrated into conveyor and PLC reject hardware, cutting inspection time by 99.5\%, and loom-side fabric structural-defect detection using a microscopic camera array, cutting operator headcount by 50\%. Both required the same discipline a medical CV system demands -- getting an algorithm from a lab result to a piece of hardware making real-time decisions on a live production line.

\vspace{0.5em}

\noindent My PhD at the University of Trento sharpened that discipline further. I designed frozen-protocol benchmark suites and ablation studies to quantify accuracy, latency, and CPU trade-offs for real-time systems under hard constraints, and, as a Visiting Researcher at the University of Visual and Performing Arts in Colombo, ran a human-centred evaluation study combining quantitative and qualitative methods to assess how end users actually perceived a deployed detection system. I see both as directly transferable to validating a computer-vision system intended for clinical use, where a benchmark's numbers matter only insofar as they hold up under real-world conditions and real-world users. My engineering background is equally hands-on: C++17 for real-time inference engines, and embedded deployment on ARM-class hardware including Raspberry Pi, disciplines I have carried from apparel-manufacturing production lines through to my own current open-source work.

\vspace{0.5em}

\noindent I want to be direct about where my background does not yet reach: I have no experience with medical-device regulatory processes such as FDA or Health Canada device approval, and no surgical-domain-specific computer-vision experience. What I bring instead is a decade of production CV discipline that has already crossed from industrial inspection into one real clinical deployment, plus a research background built specifically around rigorously validating systems before they are trusted in the field. I see the regulatory and surgical-domain specifics as learnable on top of that foundation, not as a substitute for it.

\vspace{0.5em}

\noindent I am a Canadian Permanent Resident, eligible to work in Canada without sponsorship, based in Toronto and open to relocating to Burnaby/Vancouver, and available immediately. I would welcome the chance to discuss how my background in production computer vision and evaluation rigor could contribute to Stryker's work.

\vspace{1em}

\noindent Sincerely, \\
Nishal Stanislaus Silva

\end{document}
